\documentclass[11pt,a4paper]{article}

%====================================================
% ENCODAGE ET LANGUE
%====================================================
\usepackage[utf8]{inputenc}
\usepackage[T1]{fontenc}
\usepackage[french]{babel}
\usepackage{lmodern}

%====================================================
% MISE EN PAGE
%====================================================
\usepackage[a4paper,margin=2cm]{geometry}
\usepackage{fancyhdr}
\usepackage{lastpage}
\usepackage{titlesec}
\usepackage{setspace}
\usepackage{enumitem}
\usepackage{array}
\usepackage{tabularx}
\usepackage{multicol}
\usepackage{booktabs}
\usepackage{hyperref}

%====================================================
% MATHS ET PHYSIQUE
%====================================================
\usepackage{amsmath,amssymb,amsthm}
\usepackage{siunitx}
\sisetup{
  locale = FR,
  detect-all,
  per-mode=symbol,
  separate-uncertainty=true
}

%====================================================
% FIGURES ET ENCADRÉS
%====================================================
\usepackage{tikz}
\usetikzlibrary{arrows.meta,calc,decorations.pathmorphing,patterns}
\usepackage[most]{tcolorbox}

%====================================================
% COULEURS
%====================================================
\definecolor{bleu}{RGB}{25,70,130}
\definecolor{bleuclair}{RGB}{235,243,255}
\definecolor{vert}{RGB}{35,120,70}
\definecolor{vertclair}{RGB}{235,250,240}
\definecolor{orange}{RGB}{190,100,20}
\definecolor{orangeclair}{RGB}{255,246,232}
\definecolor{rouge}{RGB}{170,40,40}
\definecolor{rougeclair}{RGB}{255,238,238}
\definecolor{violet}{RGB}{100,60,150}
\definecolor{violetclair}{RGB}{245,240,255}
\definecolor{grisclair}{RGB}{245,245,245}

%====================================================
% EN-TÊTE
%====================================================
\pagestyle{fancy}
\fancyhf{}
\fancyhead[L]{Première spécialité physique-chimie}
\fancyhead[C]{TD -- Ondes et signaux}
\fancyhead[R]{Page \thepage/\pageref{LastPage}}
\fancyfoot[C]{\small Document de travail -- Exercices progressifs et corrigés détaillés}
\setlength{\headheight}{14pt}

%====================================================
% TITRES
%====================================================
\titleformat{\section}
{\Large\bfseries\color{bleu}}
{\thesection.}
{0.5em}
{}

\titleformat{\subsection}
{\large\bfseries\color{violet}}
{\thesubsection.}
{0.5em}
{}

\titleformat{\subsubsection}
{\bfseries\color{orange}}
{\thesubsubsection.}
{0.5em}
{}

%====================================================
% ENCADRÉS
%====================================================
\tcbset{
  enhanced,
  breakable,
  arc=2mm,
  boxrule=0.7pt,
  before skip=8pt,
  after skip=8pt
}

\newtcolorbox{methodbox}{
  colback=bleuclair,
  colframe=bleu,
  title=\textbf{Méthode},
  fonttitle=\bfseries
}

\newtcolorbox{retenirbox}{
  colback=vertclair,
  colframe=vert,
  title=\textbf{À retenir},
  fonttitle=\bfseries
}

\newtcolorbox{erreurbox}{
  colback=rougeclair,
  colframe=rouge,
  title=\textbf{Erreur classique},
  fonttitle=\bfseries
}

\newtcolorbox{vigilancebox}{
  colback=orangeclair,
  colframe=orange,
  title=\textbf{Point de vigilance},
  fonttitle=\bfseries
}

\newtcolorbox{correctionbox}{
  colback=grisclair,
  colframe=black!65,
  title=\textbf{Corrigé détaillé},
  fonttitle=\bfseries
}

\newtcolorbox{exobox}{
  colback=white,
  colframe=bleu,
  fonttitle=\bfseries,
  title=\thetcbcounter,
  enhanced,
  breakable
}

%====================================================
% COMMANDES
%====================================================
\newcounter{exercice}
\newcommand{\etoileA}{$\star\circ\circ\circ\circ$}
\newcommand{\etoileB}{$\star\star\circ\circ\circ$}
\newcommand{\etoileC}{$\star\star\star\circ\circ$}
\newcommand{\etoileD}{$\star\star\star\star\circ$}
\newcommand{\etoileE}{$\star\star\star\star\star$}

\newcommand{\exo}[3]{
  \refstepcounter{exercice}
  \begin{tcolorbox}[
    colback=white,
    colframe=bleu,
    title=\textbf{Exercice \theexercice{} -- #1 -- #2},
    breakable
  ]
  #3
  \end{tcolorbox}
}

\newcommand{\corr}[2]{
  \begin{correctionbox}
  \textbf{Exercice #1.} #2
  \end{correctionbox}
}

\newcommand{\dd}{\mathrm{d}}
\newcommand{\eV}{\mathrm{eV}}
\newcommand{\Hz}{\mathrm{Hz}}

%====================================================
% DOCUMENT
%====================================================
\begin{document}

%====================================================
% PAGE DE TITRE
%====================================================
\begin{titlepage}
\begin{center}
\vspace*{1cm}

{\Huge\bfseries TD complet}\\[0.3cm]
{\Huge\bfseries Ondes et signaux}\\[0.5cm]

{\Large Première générale -- Spécialité physique-chimie}\\[1cm]

\begin{tikzpicture}
\draw[very thick, bleu, domain=0:10, samples=300] plot(\x,{0.7*sin(180*\x)});
\draw[-{Latex}, thick] (-0.2,0) -- (10.6,0) node[right] {$x$};
\draw[-{Latex}, thick] (0,-1.1) -- (0,1.2) node[above] {$y$};
\end{tikzpicture}

\vfill

\begin{tcolorbox}[colback=bleuclair,colframe=bleu,width=0.9\textwidth,title=\textbf{Objectif du TD}]
Ce TD a pour objectif de travailler de manière progressive toutes les notions du thème
\textbf{Ondes et signaux} du programme de première spécialité physique-chimie :
ondes mécaniques progressives, célérité, retard, ondes périodiques, longueur d'onde,
lentilles minces convergentes, grandissement, couleurs, domaines des ondes électromagnétiques,
modèle du photon et spectres d'émission ou d'absorption.
\end{tcolorbox}

\vfill

\begin{tabular}{rl}
\textbf{Durée indicative :} & plusieurs séances de TD ou travail progressif annuel.\\
\textbf{Nombre d'exercices :} & 32 exercices progressifs.\\
\textbf{Difficulté :} & de $\star\circ\circ\circ\circ$ à $\star\star\star\star\star$.\\
\textbf{Corrigés :} & complets et détaillés en fin de document.\\
\end{tabular}

\vfill

{\large \today}

\end{center}
\end{titlepage}

\tableofcontents
\newpage

%====================================================
\section{Rappels essentiels de méthode}

\subsection{Ondes mécaniques}

\begin{retenirbox}
Une \textbf{onde mécanique progressive} correspond à la propagation d'une perturbation dans un milieu matériel, sans transport global de matière, mais avec transport d'énergie.

La célérité $v$ d'une onde est définie par :
\[
v=\frac{d}{\Delta t},
\]
où $d$ est la distance parcourue par la perturbation et $\Delta t$ la durée de propagation.
\end{retenirbox}

\begin{methodbox}
Pour exploiter un retard temporel :
\[
\tau = \frac{d}{v}.
\]
Si deux capteurs séparés d'une distance $d$ détectent la même perturbation avec un retard $\tau$, alors :
\[
v=\frac{d}{\tau}.
\]
\end{methodbox}

\begin{retenirbox}
Pour une onde périodique sinusoïdale :
\[
v=\lambda f=\frac{\lambda}{T}.
\]
Avec :
\begin{itemize}
  \item $v$ : célérité de l'onde en \si{m.s^{-1}} ;
  \item $\lambda$ : longueur d'onde en \si{m} ;
  \item $f$ : fréquence en \si{Hz} ;
  \item $T$ : période temporelle en \si{s}.
\end{itemize}
\end{retenirbox}

\begin{erreurbox}
Ne pas confondre :
\begin{itemize}
  \item la période temporelle $T$, mesurée sur un graphe en fonction du temps ;
  \item la longueur d'onde $\lambda$, mesurée sur un graphe en fonction de la position.
\end{itemize}
\end{erreurbox}

\subsection{Lentilles minces convergentes}

\begin{retenirbox}
Pour une lentille mince convergente de centre optique $O$, de foyer image $F'$ et de distance focale $f'=\overline{OF'}$, la relation de conjugaison est :
\[
\frac{1}{\overline{OA'}}-\frac{1}{\overline{OA}}=\frac{1}{f'}.
\]
Le grandissement vaut :
\[
\gamma=\frac{\overline{A'B'}}{\overline{AB}}=\frac{\overline{OA'}}{\overline{OA}}.
\]
\end{retenirbox}

\begin{vigilancebox}
Les grandeurs $\overline{OA}$, $\overline{OA'}$, $\overline{AB}$ et $\overline{A'B'}$ sont des grandeurs algébriques. Leur signe porte une information physique : position de l'objet, position de l'image, image droite ou renversée.
\end{vigilancebox}

\subsection{Couleurs et modèles de la lumière}

\begin{retenirbox}
Pour une onde électromagnétique dans le vide ou dans l'air :
\[
c=\lambda \nu,
\]
où $c\simeq \SI{3.00e8}{m.s^{-1}}$.

L'énergie d'un photon est :
\[
E=h\nu=\frac{hc}{\lambda},
\]
avec $h=\SI{6.63e-34}{J.s}$.
\end{retenirbox}

\begin{methodbox}
Pour passer des joules aux électronvolts :
\[
\SI{1}{eV}=\SI{1.60e-19}{J}.
\]
Donc :
\[
E(\eV)=\frac{E(\si{J})}{\SI{1.60e-19}{J}}.
\]
\end{methodbox}

\newpage

%====================================================
\section{Exercices}

%====================================================
\subsection{Niveau 1 : applications directes \texorpdfstring{$\star\circ\circ\circ\circ$}{*}}

\exo{\etoileA}{Identifier une onde mécanique}{
On considère les phénomènes suivants :
\begin{enumerate}[label=\alph*)]
  \item une vague se propageant à la surface de l'eau ;
  \item un signal sonore se propageant dans l'air ;
  \item la lumière émise par une lampe ;
  \item une secousse sismique se propageant dans le sol ;
  \item une impulsion envoyée le long d'une corde tendue.
\end{enumerate}

\begin{enumerate}
  \item Parmi ces phénomènes, lesquels correspondent à des ondes mécaniques ?
  \item Pour chaque onde mécanique, préciser le milieu matériel de propagation.
  \item Expliquer pourquoi la lumière n'est pas une onde mécanique.
\end{enumerate}
}

\exo{\etoileA}{Célérité d'une impulsion sur une corde}{
Une impulsion créée sur une corde parcourt une distance de \SI{2.40}{m} en \SI{0.80}{s}.

\begin{enumerate}
  \item Rappeler la relation entre célérité, distance et durée de propagation.
  \item Calculer la célérité de l'impulsion.
  \item Interpréter le résultat obtenu par une phrase.
\end{enumerate}
}

\exo{\etoileA}{Retard d'une onde sonore}{
Un son se propage dans l'air à la célérité $v=\SI{340}{m.s^{-1}}$.
Un observateur se trouve à \SI{170}{m} d'une source sonore.

\begin{enumerate}
  \item Calculer la durée nécessaire au son pour atteindre l'observateur.
  \item Cette durée est appelée retard. Expliquer pourquoi.
\end{enumerate}
}

\exo{\etoileA}{Période et fréquence d'un signal sonore}{
Un diapason émet un son de fréquence $f=\SI{440}{Hz}$.

\begin{enumerate}
  \item Rappeler la relation entre la période $T$ et la fréquence $f$.
  \item Calculer la période du signal sonore.
  \item Donner le résultat en millisecondes.
\end{enumerate}
}

\exo{\etoileA}{Longueur d'onde d'un son}{
Un son de fréquence \SI{680}{Hz} se propage dans l'air à la célérité \SI{340}{m.s^{-1}}.

\begin{enumerate}
  \item Rappeler la relation entre célérité, longueur d'onde et fréquence.
  \item Calculer la longueur d'onde du son.
  \item Comment évolue la longueur d'onde si la fréquence augmente dans le même milieu ?
\end{enumerate}
}

\exo{\etoileA}{Lentille convergente et foyer}{
Une lentille mince convergente possède une distance focale $f'=\SI{10.0}{cm}$.

\begin{enumerate}
  \item Que représente la distance focale ?
  \item Où se situe le foyer image $F'$ par rapport au centre optique $O$ ?
  \item Calculer la vergence $C$ de la lentille en dioptries, sachant que $C=1/f'$ avec $f'$ en mètres.
\end{enumerate}
}

\exo{\etoileA}{Fréquence d'une onde électromagnétique}{
Une onde électromagnétique a une longueur d'onde dans le vide $\lambda=\SI{600}{nm}$.

\begin{enumerate}
  \item Convertir cette longueur d'onde en mètres.
  \item Calculer sa fréquence.
  \item Cette radiation appartient-elle au domaine visible ?
\end{enumerate}

On prendra $c=\SI{3.00e8}{m.s^{-1}}$.
}

\exo{\etoileA}{Énergie d'un photon}{
Un photon possède une fréquence $\nu=\SI{5.00e14}{Hz}$.

\begin{enumerate}
  \item Rappeler l'expression de l'énergie d'un photon.
  \item Calculer son énergie en joules.
  \item Convertir cette énergie en électronvolts.
\end{enumerate}

Données : $h=\SI{6.63e-34}{J.s}$ et $\SI{1}{eV}=\SI{1.60e-19}{J}$.
}

%====================================================
\subsection{Niveau 2 : exercices standards \texorpdfstring{$\star\star\circ\circ\circ$}{**}}

\exo{\etoileB}{Deux microphones et mesure de célérité}{
Deux microphones $M_1$ et $M_2$ sont alignés avec une source sonore. La distance entre les deux microphones est $d=\SI{1.20}{m}$. Le signal est reçu par $M_2$ avec un retard $\tau=\SI{3.5}{ms}$ par rapport à $M_1$.

\begin{enumerate}
  \item Faire un schéma simple de la situation.
  \item Exprimer la célérité du son en fonction de $d$ et $\tau$.
  \item Calculer la célérité du son.
  \item Comparer à la valeur usuelle \SI{340}{m.s^{-1}} et commenter.
\end{enumerate}
}

\exo{\etoileB}{Onde périodique sur une corde}{
Une onde sinusoïdale se propage sur une corde. On mesure une distance de \SI{1.80}{m} entre trois crêtes consécutives.

\begin{enumerate}
  \item Combien de longueurs d'onde sont contenues entre trois crêtes consécutives ?
  \item Déterminer la longueur d'onde $\lambda$.
  \item La fréquence de la source est \SI{12}{Hz}. Calculer la célérité de l'onde.
\end{enumerate}
}

\exo{\etoileB}{Graphe temporel d'un signal sonore}{
Un microphone enregistre un signal périodique. On observe que \SI{5}{} périodes occupent une durée totale de \SI{11.5}{ms}.

\begin{enumerate}
  \item Calculer la période $T$ du signal.
  \item En déduire sa fréquence.
  \item Ce son est-il audible par l'oreille humaine ? On rappelle que le domaine audible s'étend environ de \SI{20}{Hz} à \SI{20}{kHz}.
\end{enumerate}
}

\exo{\etoileB}{Écho dans une grotte}{
Un spéléologue émet un bref son face à une paroi. Il reçoit l'écho \SI{0.62}{s} plus tard.

\begin{enumerate}
  \item Expliquer pourquoi le son parcourt deux fois la distance entre le spéléologue et la paroi.
  \item Calculer la distance à la paroi.
  \item Que deviendrait le retard si la distance à la paroi était doublée ?
\end{enumerate}

On prendra $v_{\text{son}}=\SI{340}{m.s^{-1}}$.
}

\exo{\etoileB}{Image par une lentille convergente}{
Un objet $AB$ de hauteur \SI{2.0}{cm} est placé à \SI{30.0}{cm} devant une lentille convergente de distance focale $f'=\SI{10.0}{cm}$.

On adopte la convention :
\[
\overline{OA}=-\SI{30.0}{cm}, \qquad f'=+\SI{10.0}{cm}.
\]

\begin{enumerate}
  \item Utiliser la relation de conjugaison pour déterminer $\overline{OA'}$.
  \item Calculer le grandissement $\gamma$.
  \item En déduire la taille de l'image.
  \item L'image est-elle droite ou renversée ?
\end{enumerate}
}

\exo{\etoileB}{Couleur d'un objet éclairé}{
Une pomme rouge est éclairée successivement par :
\begin{enumerate}[label=\alph*)]
  \item une lumière blanche ;
  \item une lumière rouge ;
  \item une lumière verte ;
  \item une lumière cyan.
\end{enumerate}

On suppose que la pomme diffuse principalement le rouge et absorbe les autres couleurs.

\begin{enumerate}
  \item Indiquer la couleur perçue dans chaque situation.
  \item Expliquer le raisonnement en utilisant les mots absorption et diffusion.
\end{enumerate}
}

\exo{\etoileB}{Filtre coloré}{
Un filtre rouge est placé devant une source de lumière blanche.

\begin{enumerate}
  \item Quelles radiations sont principalement transmises ?
  \item Quelles radiations sont principalement absorbées ?
  \item Quelle couleur observe-t-on après le filtre ?
  \item Que se passe-t-il si l'on place ensuite un filtre vert derrière le filtre rouge ?
\end{enumerate}
}

\exo{\etoileB}{Domaine spectral d'une onde}{
On considère les ondes électromagnétiques suivantes :

\begin{center}
\begin{tabular}{|c|c|}
\hline
Onde & Longueur d'onde dans le vide\\
\hline
A & \SI{700}{nm}\\
B & \SI{12}{cm}\\
C & \SI{2.0}{nm}\\
D & \SI{1.0}{mm}\\
\hline
\end{tabular}
\end{center}

\begin{enumerate}
  \item Convertir toutes les longueurs d'onde en mètres.
  \item Identifier approximativement le domaine spectral de chaque onde.
  \item Calculer la fréquence de l'onde A.
\end{enumerate}
}

%====================================================
\subsection{Niveau 3 : exercices intermédiaires \texorpdfstring{$\star\star\star\circ\circ$}{***}}

\exo{\etoileC}{Localisation d'un impact sur une barre}{
Une longue barre métallique est équipée de deux capteurs $C_1$ et $C_2$ séparés de \SI{2.00}{m}. Un choc est donné entre les deux capteurs. L'onde mécanique se propage dans la barre à la célérité $v=\SI{5000}{m.s^{-1}}$.

Le signal arrive sur $C_1$ avec \SI{0.20}{ms} d'avance sur $C_2$.

\begin{enumerate}
  \item Noter $x$ la distance entre le point d'impact et $C_1$. Exprimer la distance entre le point d'impact et $C_2$.
  \item Écrire les durées de propagation vers $C_1$ et vers $C_2$.
  \item Utiliser le retard mesuré pour déterminer $x$.
  \item Conclure sur la position de l'impact.
\end{enumerate}
}

\exo{\etoileC}{Onde ultrasonore et télémètre}{
Un télémètre à ultrasons émet une salve ultrasonore vers un obstacle. L'écho revient après une durée $\Delta t=\SI{18.0}{ms}$. La célérité du son dans l'air vaut \SI{340}{m.s^{-1}}.

\begin{enumerate}
  \item Expliquer pourquoi la distance à l'obstacle n'est pas égale à $v\Delta t$.
  \item Calculer la distance à l'obstacle.
  \item Le télémètre est-il adapté pour mesurer une distance de \SI{20}{m} si l'acquisition ne peut pas dépasser \SI{100}{ms} ?
\end{enumerate}
}

\exo{\etoileC}{Onde sinusoïdale : lecture spatiale et temporelle}{
Une onde sinusoïdale se propage vers les $x$ croissants. À un instant donné, on observe que deux points consécutifs dans le même état vibratoire sont séparés de \SI{0.40}{m}. À une position fixe, la durée entre deux maxima successifs est \SI{0.050}{s}.

\begin{enumerate}
  \item Identifier la longueur d'onde.
  \item Identifier la période.
  \item Calculer la fréquence.
  \item Calculer la célérité.
  \item Écrire une phrase distinguant périodicité spatiale et périodicité temporelle.
\end{enumerate}
}

\exo{\etoileC}{Influence du milieu sur la célérité}{
Une même onde sonore de fréquence \SI{1000}{Hz} se propage :
\begin{itemize}
  \item dans l'air à \SI{340}{m.s^{-1}} ;
  \item dans l'eau à \SI{1500}{m.s^{-1}} ;
  \item dans l'acier à \SI{5000}{m.s^{-1}}.
\end{itemize}

\begin{enumerate}
  \item Calculer la longueur d'onde dans chaque milieu.
  \item La fréquence change-t-elle lorsque l'onde passe d'un milieu à un autre ?
  \item Quelle grandeur dépend du milieu ?
  \item Commenter l'idée : « la longueur d'onde dépend à la fois de la source et du milieu ».
\end{enumerate}
}

\exo{\etoileC}{Mise au point photographique}{
Un appareil photo est modélisé par une lentille convergente. Un objet très éloigné forme une image nette sur un capteur placé à \SI{50.0}{mm} de la lentille.

\begin{enumerate}
  \item En déduire la distance focale de la lentille.
  \item Un objet est maintenant placé à \SI{2.00}{m} devant la lentille. Déterminer la nouvelle position de l'image.
  \item Le capteur doit-il se rapprocher ou s'éloigner de la lentille pour retrouver une image nette ?
\end{enumerate}
}

\exo{\etoileC}{Image virtuelle avec une lentille convergente}{
Un objet est placé à \SI{5.0}{cm} devant une lentille convergente de distance focale \SI{10.0}{cm}. On prendra :
\[
\overline{OA}=-\SI{5.0}{cm}, \qquad f'=+\SI{10.0}{cm}.
\]

\begin{enumerate}
  \item Déterminer $\overline{OA'}$.
  \item L'image est-elle réelle ou virtuelle ?
  \item Calculer le grandissement.
  \item L'image est-elle droite ou renversée ?
  \item Cette situation peut-elle servir de modèle simple pour une loupe ?
\end{enumerate}
}

\exo{\etoileC}{Synthèse additive et écran}{
Un écran de téléphone utilise trois sous-pixels : rouge, vert et bleu. On suppose que chaque sous-pixel peut être allumé ou éteint.

\begin{enumerate}
  \item Quelle couleur obtient-on si seuls les sous-pixels rouge et vert sont allumés ?
  \item Quelle couleur obtient-on si les trois sous-pixels sont allumés avec la même intensité ?
  \item Quelle couleur obtient-on si rouge et bleu sont allumés ?
  \item Pourquoi parle-t-on de synthèse additive ?
  \item Pourquoi un écran noir correspond-il à l'absence d'émission lumineuse ?
\end{enumerate}
}

\exo{\etoileC}{Diagramme de niveaux d'énergie}{
Un atome possède deux niveaux d'énergie :
\[
E_1=-\SI{5.40}{eV}, \qquad E_2=-\SI{3.10}{eV}.
\]

\begin{enumerate}
  \item Calculer l'énergie du photon émis lors de la transition de $E_2$ vers $E_1$.
  \item Convertir cette énergie en joules.
  \item Calculer la fréquence du photon.
  \item Calculer sa longueur d'onde dans le vide.
  \item Dans quel domaine spectral se situe cette radiation ?
\end{enumerate}

Données :
\[
h=\SI{6.63e-34}{J.s}, \quad c=\SI{3.00e8}{m.s^{-1}}, \quad \SI{1}{eV}=\SI{1.60e-19}{J}.
\]
}

%====================================================
\subsection{Niveau 4 : exercices difficiles \texorpdfstring{$\star\star\star\star\circ$}{****}}

\exo{\etoileD}{Mini-sismologie : localisation d'un épicentre}{
Lors d'un séisme, deux types d'ondes sont détectées par une station :
\begin{itemize}
  \item les ondes P, rapides, de célérité $v_P=\SI{6.0}{km.s^{-1}}$ ;
  \item les ondes S, plus lentes, de célérité $v_S=\SI{3.5}{km.s^{-1}}$.
\end{itemize}

Dans une station, les ondes S arrivent \SI{42}{s} après les ondes P.

\begin{enumerate}
  \item Si $D$ est la distance entre l'épicentre et la station, exprimer les durées $t_P$ et $t_S$.
  \item Écrire l'expression du retard $\Delta t=t_S-t_P$.
  \item Déterminer la distance $D$.
  \item Expliquer pourquoi une seule station ne suffit pas en général à localiser précisément l'épicentre sur une carte.
\end{enumerate}
}

\exo{\etoileD}{Corde vibrante et modèle sinusoïdal}{
Une corde est parcourue par une onde sinusoïdale décrite par :
\[
y(x,t)=A\cos\left(2\pi\left(\frac{t}{T}-\frac{x}{\lambda}\right)\right),
\]
avec $A=\SI{2.0}{cm}$, $T=\SI{0.10}{s}$ et $\lambda=\SI{0.50}{m}$.

\begin{enumerate}
  \item Identifier l'amplitude, la période et la longueur d'onde.
  \item Calculer la fréquence.
  \item Calculer la célérité de l'onde.
  \item Déterminer $y(0,0)$.
  \item Déterminer $y(\lambda/2,0)$.
  \item Interpréter physiquement le signe négatif devant $x/\lambda$.
\end{enumerate}
}

\exo{\etoileD}{Lentille à focale variable}{
Une lentille à focale variable est utilisée pour former sur un capteur une image nette d'un objet placé à \SI{40.0}{cm} devant la lentille. Le capteur est fixe à \SI{20.0}{cm} derrière la lentille.

\begin{enumerate}
  \item Donner les valeurs algébriques de $\overline{OA}$ et $\overline{OA'}$.
  \item Déterminer la distance focale nécessaire.
  \item Calculer la vergence correspondante.
  \item L'image est-elle droite ou renversée ?
  \item Si l'objet se rapproche, faut-il augmenter ou diminuer la distance focale pour garder le capteur fixe ? Justifier qualitativement.
\end{enumerate}
}

\exo{\etoileD}{Filtres en photographie astronomique}{
Une nébuleuse émet principalement trois radiations visibles :
\[
\lambda_1=\SI{486}{nm},\qquad \lambda_2=\SI{656}{nm},\qquad \lambda_3=\SI{501}{nm}.
\]
Un photographe utilise successivement :
\begin{itemize}
  \item un filtre rouge transmettant surtout les longueurs d'onde supérieures à \SI{600}{nm} ;
  \item un filtre bleu transmettant surtout les longueurs d'onde entre \SI{430}{nm} et \SI{500}{nm} ;
  \item un filtre vert transmettant surtout les longueurs d'onde entre \SI{500}{nm} et \SI{570}{nm}.
\end{itemize}

\begin{enumerate}
  \item Pour chaque filtre, indiquer quelle radiation est principalement transmise.
  \item Calculer la fréquence de la radiation $\lambda_2$.
  \item Expliquer pourquoi un filtre permet d'isoler certaines informations sur la composition ou l'état d'un objet astronomique.
\end{enumerate}
}

\exo{\etoileD}{Spectre et transition atomique}{
Une raie spectrale observée dans le visible possède une longueur d'onde $\lambda=\SI{589}{nm}$.

\begin{enumerate}
  \item Calculer l'énergie du photon correspondant en joules.
  \item Convertir cette énergie en électronvolts.
  \item Cette raie peut-elle correspondre à une transition entre deux niveaux séparés de \SI{2.11}{eV} ?
  \item S'agit-il d'un spectre d'émission ou d'absorption si la raie apparaît brillante sur fond noir ?
  \item Interpréter qualitativement ce que signifie la quantification des niveaux d'énergie.
\end{enumerate}
}

%====================================================
\subsection{Niveau 5 : exercices ambitieux \texorpdfstring{$\star\star\star\star\star$}{*****}}

\exo{\etoileE}{Acoustique d'une salle de concert}{
Dans une salle de concert, un musicien produit un bref claquement. Un microphone reçoit d'abord le son direct, puis un premier écho dû à la réflexion sur le mur du fond. L'écho arrive \SI{0.12}{s} après le son direct.

Le microphone est placé à \SI{8.0}{m} du musicien. Le mur du fond est situé derrière le microphone, sur le même axe que le musicien et le microphone. On note $L$ la distance entre le musicien et le mur.

\begin{enumerate}
  \item Représenter la situation sur un schéma.
  \item Exprimer la distance parcourue par le son direct.
  \item Exprimer la distance parcourue par le son réfléchi sur le mur.
  \item Utiliser le retard pour déterminer $L$.
  \item Expliquer pourquoi les échos trop tardifs peuvent nuire à la qualité acoustique d'une salle.
\end{enumerate}

On prendra $v_{\text{son}}=\SI{340}{m.s^{-1}}$.
}

\exo{\etoileE}{Imagerie ultrasonore médicale}{
Une sonde échographique émet des ultrasons de fréquence \SI{5.0}{MHz} dans les tissus biologiques. La célérité moyenne des ultrasons dans ces tissus est \SI{1540}{m.s^{-1}}.

Un écho provenant d'une interface entre deux tissus est détecté \SI{52}{\micro s} après l'émission.

\begin{enumerate}
  \item Calculer la longueur d'onde des ultrasons dans les tissus.
  \item Calculer la profondeur de l'interface responsable de l'écho.
  \item Expliquer pourquoi les ultrasons, et non les sons audibles, sont utilisés pour former des images médicales précises.
  \item Si la fréquence augmente, que devient la longueur d'onde ? Quel peut être l'intérêt pour l'imagerie ?
  \item Quel compromis physique peut-on imaginer entre précision de l'image et pénétration dans les tissus ?
\end{enumerate}
}

\exo{\etoileE}{Synthèse finale : onde, lentille et photon}{
Un dispositif scientifique utilise :
\begin{itemize}
  \item une source lumineuse monochromatique de longueur d'onde $\lambda=\SI{532}{nm}$ ;
  \item une lentille convergente de distance focale $f'=\SI{12.0}{cm}$ ;
  \item un objet lumineux de hauteur \SI{1.5}{cm} placé à \SI{36.0}{cm} devant la lentille.
\end{itemize}

\begin{enumerate}
  \item Calculer la fréquence de la lumière émise.
  \item Calculer l'énergie d'un photon en joules puis en électronvolts.
  \item Déterminer la position de l'image formée par la lentille.
  \item Calculer le grandissement et la taille de l'image.
  \item L'image est-elle réelle ou virtuelle ? droite ou renversée ?
  \item Expliquer pourquoi ce problème mobilise à la fois le modèle géométrique, le modèle ondulatoire et le modèle particulaire de la lumière.
\end{enumerate}
}

\newpage

%====================================================
\section{Exercices bilan supplémentaires}

\exo{\etoileD}{Bilan 1 -- Concert, microphones et couleur}{
Lors d'un concert, deux microphones séparés de \SI{2.04}{m} reçoivent le même son avec un retard de \SI{6.0}{ms}. Une caméra filme la scène. Elle utilise un capteur sensible au rouge, au vert et au bleu.

\begin{enumerate}
  \item Calculer la célérité du son mesurée par les microphones.
  \item Comparer à la valeur usuelle dans l'air.
  \item Un son de fréquence \SI{440}{Hz} est joué. Calculer sa longueur d'onde dans l'air.
  \item Expliquer comment un écran de caméra restitue une couleur jaune à partir de sous-pixels.
  \item Pourquoi ne faut-il pas confondre onde sonore et onde lumineuse ?
\end{enumerate}
}

\exo{\etoileE}{Bilan 2 -- Observation d'une étoile}{
Une étoile émet une raie lumineuse de longueur d'onde \SI{656}{nm}. On utilise une lentille convergente de distance focale \SI{20.0}{cm} pour former son image sur un capteur.

L'étoile étant très éloignée, on peut considérer que ses rayons arrivent parallèles sur la lentille.

\begin{enumerate}
  \item Où doit-on placer le capteur pour obtenir une image nette ?
  \item Calculer la fréquence de la radiation observée.
  \item Calculer l'énergie d'un photon correspondant.
  \item Convertir cette énergie en électronvolts.
  \item Si cette raie correspond à une transition atomique, quelle est la différence d'énergie entre les deux niveaux ?
  \item La raie observée est rouge. Justifier à partir de la longueur d'onde.
\end{enumerate}
}

\newpage

%====================================================
\section{Corrigés détaillés}

\corr{1}{
\begin{enumerate}
  \item Les ondes mécaniques sont : la vague à la surface de l'eau, le son dans l'air, la secousse sismique dans le sol et l'impulsion le long d'une corde.
  \item Les milieux matériels sont respectivement : l'eau, l'air, le sol et la corde.
  \item La lumière peut se propager dans le vide. Elle n'a donc pas besoin d'un milieu matériel : ce n'est pas une onde mécanique mais une onde électromagnétique.
\end{enumerate}
}

\corr{2}{
\begin{enumerate}
  \item La relation est :
  \[
  v=\frac{d}{\Delta t}.
  \]
  \item On remplace :
  \[
  v=\frac{\SI{2.40}{m}}{\SI{0.80}{s}}=\SI{3.0}{m.s^{-1}}.
  \]
  \item L'impulsion se propage le long de la corde à la vitesse de \SI{3.0}{m.s^{-1}}.
\end{enumerate}
}

\corr{3}{
\begin{enumerate}
  \item On utilise :
  \[
  \Delta t=\frac{d}{v}.
  \]
  Donc :
  \[
  \Delta t=\frac{170}{340}=\SI{0.50}{s}.
  \]
  \item Le son n'est pas reçu instantanément : il met \SI{0.50}{s} pour aller de la source à l'observateur. Cette durée est donc un retard de réception.
\end{enumerate}
}

\corr{4}{
\begin{enumerate}
  \item On a :
  \[
  T=\frac{1}{f}.
  \]
  \item Donc :
  \[
  T=\frac{1}{440}=\SI{2.27e-3}{s}.
  \]
  \item En millisecondes :
  \[
  T=\SI{2.27}{ms}.
  \]
\end{enumerate}
}

\corr{5}{
\begin{enumerate}
  \item La relation est :
  \[
  v=\lambda f.
  \]
  \item Donc :
  \[
  \lambda=\frac{v}{f}=\frac{340}{680}=\SI{0.50}{m}.
  \]
  \item Dans un même milieu, $v$ est fixée. Si $f$ augmente, alors $\lambda=v/f$ diminue.
\end{enumerate}
}

\corr{6}{
\begin{enumerate}
  \item La distance focale est la distance algébrique entre le centre optique $O$ et le foyer image $F'$.
  \item Pour une lentille convergente, $F'$ est situé du côté de la lumière émergente, à \SI{10.0}{cm} de $O$.
  \item Il faut convertir :
  \[
  f'=\SI{10.0}{cm}=\SI{0.100}{m}.
  \]
  Alors :
  \[
  C=\frac{1}{f'}=\frac{1}{0.100}=\SI{10.0}{\delta}.
  \]
\end{enumerate}
}

\corr{7}{
\begin{enumerate}
  \item
  \[
  \lambda=\SI{600}{nm}=\SI{600e-9}{m}=\SI{6.00e-7}{m}.
  \]
  \item
  \[
  \nu=\frac{c}{\lambda}=\frac{\SI{3.00e8}{m.s^{-1}}}{\SI{6.00e-7}{m}}
  =\SI{5.00e14}{Hz}.
  \]
  \item \SI{600}{nm} appartient au domaine visible, plus précisément vers l'orange-rouge.
\end{enumerate}
}

\corr{8}{
\begin{enumerate}
  \item L'énergie d'un photon est :
  \[
  E=h\nu.
  \]
  \item
  \[
  E=\SI{6.63e-34}{J.s}\times \SI{5.00e14}{Hz}
  =\SI{3.315e-19}{J}.
  \]
  \item
  \[
  E=\frac{\SI{3.315e-19}{J}}{\SI{1.60e-19}{J/eV}}
  =\SI{2.07}{eV}.
  \]
\end{enumerate}
}

\corr{9}{
\begin{enumerate}
  \item Le son arrive d'abord sur $M_1$, puis sur $M_2$. Les microphones sont séparés de \SI{1.20}{m}.
  \item La célérité vaut :
  \[
  v=\frac{d}{\tau}.
  \]
  \item Attention à convertir :
  \[
  \tau=\SI{3.5}{ms}=\SI{3.5e-3}{s}.
  \]
  Donc :
  \[
  v=\frac{1.20}{3.5\times 10^{-3}}=\SI{343}{m.s^{-1}}.
  \]
  \item Cette valeur est très proche de \SI{340}{m.s^{-1}}. L'écart peut venir des incertitudes de mesure ou des conditions de température.
\end{enumerate}
}

\corr{10}{
\begin{enumerate}
  \item Entre trois crêtes consécutives, il y a deux intervalles crête-crête, donc deux longueurs d'onde.
  \item
  \[
  2\lambda=\SI{1.80}{m}
  \quad \Rightarrow \quad
  \lambda=\SI{0.90}{m}.
  \]
  \item
  \[
  v=\lambda f=0.90\times 12=\SI{10.8}{m.s^{-1}}.
  \]
\end{enumerate}
}

\corr{11}{
\begin{enumerate}
  \item Cinq périodes durent \SI{11.5}{ms}, donc :
  \[
  T=\frac{11.5}{5}=\SI{2.30}{ms}=\SI{2.30e-3}{s}.
  \]
  \item
  \[
  f=\frac{1}{T}=\frac{1}{2.30\times 10^{-3}}=\SI{435}{Hz}.
  \]
  \item Ce son est audible car sa fréquence est comprise entre \SI{20}{Hz} et \SI{20}{kHz}.
\end{enumerate}
}

\corr{12}{
\begin{enumerate}
  \item Le son part du spéléologue, atteint la paroi, puis revient au spéléologue. Il parcourt donc $2d$.
  \item
  \[
  2d=v\Delta t
  \quad \Rightarrow \quad
  d=\frac{v\Delta t}{2}.
  \]
  Donc :
  \[
  d=\frac{340\times 0.62}{2}=\SI{105.4}{m}.
  \]
  La paroi est à environ \SI{105}{m}.
  \item Si la distance double, le trajet aller-retour double, donc le retard double.
\end{enumerate}
}

\corr{13}{
\begin{enumerate}
  \item Relation de conjugaison :
  \[
  \frac{1}{\overline{OA'}}-\frac{1}{\overline{OA}}=\frac{1}{f'}.
  \]
  Avec $\overline{OA}=-\SI{30.0}{cm}$ et $f'=+\SI{10.0}{cm}$ :
  \[
  \frac{1}{\overline{OA'}}-\frac{1}{-30.0}=\frac{1}{10.0}.
  \]
  Donc :
  \[
  \frac{1}{\overline{OA'}}+\frac{1}{30.0}=\frac{1}{10.0}.
  \]
  Ainsi :
  \[
  \frac{1}{\overline{OA'}}=\frac{1}{10.0}-\frac{1}{30.0}
  =\frac{3-1}{30}=\frac{2}{30}=\frac{1}{15}.
  \]
  Donc :
  \[
  \overline{OA'}=+\SI{15.0}{cm}.
  \]
  \item
  \[
  \gamma=\frac{\overline{OA'}}{\overline{OA}}=\frac{15.0}{-30.0}=-0.50.
  \]
  \item
  \[
  \overline{A'B'}=\gamma \overline{AB}=-0.50\times \SI{2.0}{cm}=-\SI{1.0}{cm}.
  \]
  \item Le signe négatif indique une image renversée. Elle est réelle car $\overline{OA'}>0$.
\end{enumerate}
}

\corr{14}{
\begin{enumerate}
  \item En lumière blanche, la pomme diffuse le rouge : elle apparaît rouge.
  \item En lumière rouge, elle reçoit du rouge et le diffuse : elle apparaît rouge.
  \item En lumière verte, elle absorbe principalement la lumière reçue et ne diffuse presque rien : elle apparaît noire ou très sombre.
  \item La lumière cyan contient du vert et du bleu, mais pas de rouge. La pomme rouge ne peut pas diffuser de rouge : elle apparaît sombre.
\end{enumerate}
}

\corr{15}{
\begin{enumerate}
  \item Un filtre rouge transmet principalement le rouge.
  \item Il absorbe principalement le vert et le bleu.
  \item Après le filtre, la lumière observée est rouge.
  \item Le filtre rouge laisse passer du rouge. Le filtre vert placé ensuite absorbe le rouge. Il ne reste presque aucune lumière : on observe du noir ou une forte obscurité.
\end{enumerate}
}

\corr{16}{
\begin{enumerate}
  \item
  \[
  \SI{700}{nm}=\SI{7.00e-7}{m},\quad
  \SI{12}{cm}=\SI{1.2e-1}{m},\quad
  \SI{2.0}{nm}=\SI{2.0e-9}{m},\quad
  \SI{1.0}{mm}=\SI{1.0e-3}{m}.
  \]
  \item A est dans le visible rouge ; B est dans le domaine des ondes radio ou micro-ondes ; C est dans le domaine des rayons X ; D est dans le domaine micro-ondes ou infrarouge lointain selon les limites choisies.
  \item Pour A :
  \[
  \nu=\frac{c}{\lambda}=\frac{3.00\times 10^8}{7.00\times 10^{-7}}
  =\SI{4.29e14}{Hz}.
  \]
\end{enumerate}
}

\corr{17}{
\begin{enumerate}
  \item Si $x$ est la distance impact-$C_1$, alors la distance impact-$C_2$ est $2.00-x$.
  \item Les durées sont :
  \[
  t_1=\frac{x}{v},\qquad t_2=\frac{2.00-x}{v}.
  \]
  \item Le signal arrive sur $C_1$ avec \SI{0.20}{ms} d'avance, donc :
  \[
  t_2-t_1=\SI{0.20}{ms}=\SI{2.0e-4}{s}.
  \]
  Ainsi :
  \[
  \frac{2.00-x}{5000}-\frac{x}{5000}=2.0\times 10^{-4}.
  \]
  \[
  \frac{2.00-2x}{5000}=2.0\times 10^{-4}.
  \]
  \[
  2.00-2x=5000\times 2.0\times 10^{-4}=1.0.
  \]
  \[
  2x=1.0 \quad \Rightarrow \quad x=\SI{0.50}{m}.
  \]
  \item L'impact a eu lieu à \SI{0.50}{m} de $C_1$, donc plus près de $C_1$ que de $C_2$.
\end{enumerate}
}

\corr{18}{
\begin{enumerate}
  \item L'onde fait un aller-retour entre le télémètre et l'obstacle. La distance totale parcourue est donc $2d$.
  \item
  \[
  d=\frac{v\Delta t}{2}
  =\frac{340\times 18.0\times 10^{-3}}{2}
  =\SI{3.06}{m}.
  \]
  \item Pour \SI{20}{m}, l'aller-retour vaut \SI{40}{m}. La durée serait :
  \[
  \Delta t=\frac{40}{340}=\SI{0.118}{s}=\SI{118}{ms}.
  \]
  Comme \SI{118}{ms} dépasse \SI{100}{ms}, ce télémètre ne serait pas adapté sans modifier l'acquisition.
\end{enumerate}
}

\corr{19}{
\begin{enumerate}
  \item Deux points consécutifs dans le même état vibratoire sont séparés d'une longueur d'onde :
  \[
  \lambda=\SI{0.40}{m}.
  \]
  \item La durée entre deux maxima successifs est la période :
  \[
  T=\SI{0.050}{s}.
  \]
  \item
  \[
  f=\frac{1}{T}=\frac{1}{0.050}=\SI{20}{Hz}.
  \]
  \item
  \[
  v=\frac{\lambda}{T}=\frac{0.40}{0.050}=\SI{8.0}{m.s^{-1}}.
  \]
  \item La périodicité spatiale décrit la répétition de l'état de l'onde dans l'espace à un instant donné ; la périodicité temporelle décrit la répétition du signal au cours du temps en un point fixé.
\end{enumerate}
}

\corr{20}{
\begin{enumerate}
  \item On utilise $\lambda=v/f$.
  \[
  \lambda_{\text{air}}=\frac{340}{1000}=\SI{0.340}{m}.
  \]
  \[
  \lambda_{\text{eau}}=\frac{1500}{1000}=\SI{1.50}{m}.
  \]
  \[
  \lambda_{\text{acier}}=\frac{5000}{1000}=\SI{5.00}{m}.
  \]
  \item La fréquence est imposée par la source. Elle ne change pas lors du changement de milieu.
  \item La célérité dépend du milieu. Comme $\lambda=v/f$, la longueur d'onde dépend aussi du milieu.
  \item La fréquence dépend de la source et la célérité dépend du milieu. Comme $\lambda=v/f$, la longueur d'onde dépend donc à la fois de la source et du milieu.
\end{enumerate}
}

\corr{21}{
\begin{enumerate}
  \item Pour un objet très éloigné, l'image se forme pratiquement dans le plan focal image. Donc :
  \[
  f'=\SI{50.0}{mm}.
  \]
  \item On a :
  \[
  \overline{OA}=-\SI{2.00}{m},\qquad f'=\SI{0.0500}{m}.
  \]
  Relation de conjugaison :
  \[
  \frac{1}{\overline{OA'}}-\frac{1}{-2.00}=\frac{1}{0.0500}.
  \]
  \[
  \frac{1}{\overline{OA'}}+0.500=20.0.
  \]
  \[
  \frac{1}{\overline{OA'}}=19.5.
  \]
  \[
  \overline{OA'}=\SI{0.0513}{m}=\SI{51.3}{mm}.
  \]
  \item Le capteur doit s'éloigner légèrement de la lentille, de \SI{50.0}{mm} à environ \SI{51.3}{mm}.
\end{enumerate}
}

\corr{22}{
\begin{enumerate}
  \item Relation de conjugaison :
  \[
  \frac{1}{\overline{OA'}}-\frac{1}{\overline{OA}}=\frac{1}{f'}.
  \]
  Avec $\overline{OA}=-\SI{5.0}{cm}$ et $f'=+\SI{10.0}{cm}$ :
  \[
  \frac{1}{\overline{OA'}}+\frac{1}{5.0}=\frac{1}{10.0}.
  \]
  \[
  \frac{1}{\overline{OA'}}=\frac{1}{10.0}-\frac{1}{5.0}=-\frac{1}{10.0}.
  \]
  Donc :
  \[
  \overline{OA'}=-\SI{10.0}{cm}.
  \]
  \item L'image est virtuelle car $\overline{OA'}<0$ : elle se situe du même côté que l'objet.
  \item
  \[
  \gamma=\frac{\overline{OA'}}{\overline{OA}}=\frac{-10.0}{-5.0}=+2.0.
  \]
  \item Le grandissement est positif : l'image est droite. Elle est agrandie d'un facteur 2.
  \item Oui, cette situation modélise une loupe : l'objet est placé entre la lentille et son foyer, ce qui donne une image virtuelle, droite et agrandie.
\end{enumerate}
}

\corr{23}{
\begin{enumerate}
  \item Rouge + vert donne jaune en synthèse additive.
  \item Rouge + vert + bleu avec même intensité donne blanc.
  \item Rouge + bleu donne magenta.
  \item On parle de synthèse additive car on ajoute des lumières colorées émises.
  \item Un écran noir correspond à l'absence de lumière émise par les sous-pixels.
\end{enumerate}
}

\corr{24}{
\begin{enumerate}
  \item
  \[
  \Delta E=E_2-E_1=(-3.10)-(-5.40)=\SI{2.30}{eV}.
  \]
  \item
  \[
  \Delta E=\SI{2.30}{eV}\times \SI{1.60e-19}{J/eV}
  =\SI{3.68e-19}{J}.
  \]
  \item
  \[
  \nu=\frac{E}{h}
  =\frac{3.68\times 10^{-19}}{6.63\times 10^{-34}}
  =\SI{5.55e14}{Hz}.
  \]
  \item
  \[
  \lambda=\frac{c}{\nu}
  =\frac{3.00\times 10^8}{5.55\times 10^{14}}
  =\SI{5.41e-7}{m}
  =\SI{541}{nm}.
  \]
  \item Cette radiation appartient au domaine visible, vers le vert.
\end{enumerate}
}

\corr{25}{
\begin{enumerate}
  \item
  \[
  t_P=\frac{D}{v_P},\qquad t_S=\frac{D}{v_S}.
  \]
  \item
  \[
  \Delta t=t_S-t_P=\frac{D}{v_S}-\frac{D}{v_P}
  =D\left(\frac{1}{v_S}-\frac{1}{v_P}\right).
  \]
  \item
  \[
  D=\frac{\Delta t}{\frac{1}{v_S}-\frac{1}{v_P}}.
  \]
  Avec les vitesses en \si{km.s^{-1}} :
  \[
  D=\frac{42}{\frac{1}{3.5}-\frac{1}{6.0}}.
  \]
  \[
  \frac{1}{3.5}-\frac{1}{6.0}=0.2857-0.1667=0.1190.
  \]
  \[
  D=\frac{42}{0.1190}=\SI{353}{km}.
  \]
  \item Une station donne seulement une distance à l'épicentre. Sur une carte, cela correspond à un cercle de positions possibles. Il faut plusieurs stations pour localiser l'épicentre par intersection.
\end{enumerate}
}

\corr{26}{
\begin{enumerate}
  \item On lit :
  \[
  A=\SI{2.0}{cm},\qquad T=\SI{0.10}{s},\qquad \lambda=\SI{0.50}{m}.
  \]
  \item
  \[
  f=\frac{1}{T}=\frac{1}{0.10}=\SI{10}{Hz}.
  \]
  \item
  \[
  v=\frac{\lambda}{T}=0.50/0.10=\SI{5.0}{m.s^{-1}}.
  \]
  \item
  \[
  y(0,0)=A\cos(0)=A=\SI{2.0}{cm}.
  \]
  \item
  \[
  y(\lambda/2,0)=A\cos\left(2\pi\left(0-\frac{1}{2}\right)\right)
  =A\cos(-\pi)=-A=-\SI{2.0}{cm}.
  \]
  \item Le terme $\frac{t}{T}-\frac{x}{\lambda}$ correspond à une propagation vers les $x$ croissants.
\end{enumerate}
}

\corr{27}{
\begin{enumerate}
  \item
  \[
  \overline{OA}=-\SI{40.0}{cm},\qquad \overline{OA'}=+\SI{20.0}{cm}.
  \]
  \item Relation :
  \[
  \frac{1}{\overline{OA'}}-\frac{1}{\overline{OA}}=\frac{1}{f'}.
  \]
  \[
  \frac{1}{f'}=\frac{1}{20.0}-\frac{1}{-40.0}
  =\frac{1}{20.0}+\frac{1}{40.0}
  =\frac{3}{40.0}.
  \]
  Donc :
  \[
  f'=\frac{40.0}{3}=\SI{13.3}{cm}.
  \]
  \item En mètres :
  \[
  f'=\SI{0.133}{m}.
  \]
  Donc :
  \[
  C=\frac{1}{f'}=\SI{7.5}{\delta}.
  \]
  \item
  \[
  \gamma=\frac{20.0}{-40.0}=-0.50.
  \]
  L'image est renversée.
  \item Si l'objet se rapproche, pour garder $\overline{OA'}$ constant, la distance focale doit diminuer : la lentille doit devenir plus convergente.
\end{enumerate}
}

\corr{28}{
\begin{enumerate}
  \item Le filtre rouge transmet principalement $\lambda_2=\SI{656}{nm}$. Le filtre bleu transmet principalement $\lambda_1=\SI{486}{nm}$. Le filtre vert transmet principalement $\lambda_3=\SI{501}{nm}$.
  \item
  \[
  \nu_2=\frac{c}{\lambda_2}
  =\frac{3.00\times 10^8}{656\times 10^{-9}}
  =\SI{4.57e14}{Hz}.
  \]
  \item Les radiations émises par un gaz dépendent des transitions d'énergie possibles dans les atomes ou ions présents. En isolant une longueur d'onde, on peut obtenir des informations sur la composition chimique ou les conditions physiques de la nébuleuse.
\end{enumerate}
}

\corr{29}{
\begin{enumerate}
  \item
  \[
  E=\frac{hc}{\lambda}
  =\frac{6.63\times 10^{-34}\times 3.00\times 10^8}{589\times 10^{-9}}
  =\SI{3.38e-19}{J}.
  \]
  \item
  \[
  E=\frac{3.38\times 10^{-19}}{1.60\times 10^{-19}}
  =\SI{2.11}{eV}.
  \]
  \item Oui, car l'énergie du photon est précisément environ \SI{2.11}{eV}.
  \item Une raie brillante sur fond noir correspond à un spectre d'émission.
  \item La quantification signifie que les atomes ne peuvent posséder que certaines valeurs d'énergie. Les photons émis ou absorbés correspondent aux différences entre ces niveaux.
\end{enumerate}
}

\corr{30}{
\begin{enumerate}
  \item Schéma : musicien, microphone à \SI{8.0}{m}, mur à une distance $L$ du musicien.
  \item Le son direct parcourt :
  \[
  d_{\text{direct}}=\SI{8.0}{m}.
  \]
  \item Le son réfléchi parcourt la distance musicien-mur, puis mur-microphone :
  \[
  d_{\text{réfléchi}}=L+(L-8.0)=2L-8.0.
  \]
  \item Le retard vaut :
  \[
  \Delta t=\frac{d_{\text{réfléchi}}-d_{\text{direct}}}{v}.
  \]
  Donc :
  \[
  0.12=\frac{(2L-8.0)-8.0}{340}
  =\frac{2L-16.0}{340}.
  \]
  \[
  2L-16.0=40.8.
  \]
  \[
  2L=56.8
  \quad \Rightarrow \quad
  L=\SI{28.4}{m}.
  \]
  \item Des échos trop tardifs se superposent mal au son direct : ils peuvent rendre la musique confuse ou diminuer l'intelligibilité.
\end{enumerate}
}

\corr{31}{
\begin{enumerate}
  \item
  \[
  \lambda=\frac{v}{f}
  =\frac{1540}{5.0\times 10^6}
  =\SI{3.08e-4}{m}
  =\SI{0.308}{mm}.
  \]
  \item L'écho correspond à un aller-retour :
  \[
  d=\frac{v\Delta t}{2}
  =\frac{1540\times 52\times 10^{-6}}{2}
  =\SI{4.00e-2}{m}
  =\SI{4.0}{cm}.
  \]
  \item Les ultrasons ont des longueurs d'onde beaucoup plus petites que les sons audibles, ce qui permet de distinguer des détails plus fins.
  \item Si la fréquence augmente, $\lambda=v/f$ diminue. Une longueur d'onde plus petite peut améliorer la résolution.
  \item Les fréquences élevées donnent une meilleure résolution mais sont souvent plus atténuées dans les tissus. Il existe donc un compromis entre précision et profondeur d'exploration.
\end{enumerate}
}

\corr{32}{
\begin{enumerate}
  \item
  \[
  \nu=\frac{c}{\lambda}
  =\frac{3.00\times 10^8}{532\times 10^{-9}}
  =\SI{5.64e14}{Hz}.
  \]
  \item
  \[
  E=h\nu=6.63\times 10^{-34}\times 5.64\times 10^{14}
  =\SI{3.74e-19}{J}.
  \]
  \[
  E=\frac{3.74\times 10^{-19}}{1.60\times 10^{-19}}
  =\SI{2.34}{eV}.
  \]
  \item Avec $\overline{OA}=-\SI{36.0}{cm}$ et $f'=+\SI{12.0}{cm}$ :
  \[
  \frac{1}{\overline{OA'}}-\frac{1}{-36.0}=\frac{1}{12.0}.
  \]
  \[
  \frac{1}{\overline{OA'}}+\frac{1}{36.0}=\frac{1}{12.0}.
  \]
  \[
  \frac{1}{\overline{OA'}}=\frac{1}{12.0}-\frac{1}{36.0}
  =\frac{3-1}{36}=\frac{2}{36}=\frac{1}{18}.
  \]
  Donc :
  \[
  \overline{OA'}=+\SI{18.0}{cm}.
  \]
  \item
  \[
  \gamma=\frac{\overline{OA'}}{\overline{OA}}
  =\frac{18.0}{-36.0}=-0.50.
  \]
  \[
  \overline{A'B'}=-0.50\times \SI{1.5}{cm}
  =-\SI{0.75}{cm}.
  \]
  \item L'image est réelle car $\overline{OA'}>0$ et renversée car $\gamma<0$.
  \item Le modèle géométrique intervient avec la lentille et l'image ; le modèle ondulatoire intervient avec la fréquence et la longueur d'onde ; le modèle particulaire intervient avec l'énergie du photon.
\end{enumerate}
}

\corr{Bilan 1}{
\begin{enumerate}
  \item
  \[
  v=\frac{d}{\tau}=\frac{2.04}{6.0\times 10^{-3}}=\SI{340}{m.s^{-1}}.
  \]
  \item Cette valeur est égale à la valeur usuelle du son dans l'air à température ordinaire.
  \item
  \[
  \lambda=\frac{v}{f}=\frac{340}{440}=\SI{0.773}{m}.
  \]
  \item Sur un écran, le jaune est obtenu par synthèse additive du rouge et du vert.
  \item Une onde sonore est une onde mécanique nécessitant un milieu matériel. Une onde lumineuse est une onde électromagnétique pouvant se propager dans le vide.
\end{enumerate}
}

\corr{Bilan 2}{
\begin{enumerate}
  \item Une étoile très éloignée envoie des rayons presque parallèles. L'image se forme dans le plan focal image. Il faut placer le capteur à :
  \[
  \overline{OF'}=f'=\SI{20.0}{cm}.
  \]
  \item
  \[
  \nu=\frac{c}{\lambda}
  =\frac{3.00\times 10^8}{656\times 10^{-9}}
  =\SI{4.57e14}{Hz}.
  \]
  \item
  \[
  E=h\nu=6.63\times 10^{-34}\times 4.57\times 10^{14}
  =\SI{3.03e-19}{J}.
  \]
  \item
  \[
  E=\frac{3.03\times 10^{-19}}{1.60\times 10^{-19}}
  =\SI{1.89}{eV}.
  \]
  \item La différence d'énergie entre les deux niveaux vaut l'énergie du photon :
  \[
  \Delta E=\SI{1.89}{eV}.
  \]
  \item Une longueur d'onde de \SI{656}{nm} appartient au domaine rouge du visible.
\end{enumerate}
}

\newpage

%====================================================
\section{Fiche méthode finale}

\begin{retenirbox}
\textbf{Ondes mécaniques}
\[
v=\frac{d}{\Delta t},
\qquad
\tau=\frac{d}{v},
\qquad
v=\lambda f=\frac{\lambda}{T}.
\]
Une onde mécanique nécessite un milieu matériel. Elle transporte de l'énergie sans transport global de matière.
\end{retenirbox}

\begin{methodbox}
\textbf{Lire une onde périodique}
\begin{itemize}
  \item Sur un graphe temporel : on mesure la période $T$.
  \item Sur un graphe spatial : on mesure la longueur d'onde $\lambda$.
  \item On en déduit la célérité avec $v=\lambda/T$ ou $v=\lambda f$.
\end{itemize}
\end{methodbox}

\begin{retenirbox}
\textbf{Lentilles minces convergentes}
\[
\frac{1}{\overline{OA'}}-\frac{1}{\overline{OA}}=\frac{1}{f'},
\qquad
\gamma=\frac{\overline{A'B'}}{\overline{AB}}
=\frac{\overline{OA'}}{\overline{OA}}.
\]
Si $\overline{OA'}>0$, l'image est réelle. Si $\gamma<0$, l'image est renversée.
\end{retenirbox}

\begin{erreurbox}
\textbf{Erreurs fréquentes}
\begin{itemize}
  \item oublier de convertir les millisecondes en secondes ;
  \item confondre fréquence et période ;
  \item confondre longueur d'onde et distance parcourue ;
  \item oublier le facteur 2 dans un problème d'écho ;
  \item oublier les signes algébriques en optique ;
  \item utiliser $\lambda=c\nu$ au lieu de $c=\lambda\nu$ ;
  \item oublier de convertir les nanomètres en mètres.
\end{itemize}
\end{erreurbox}

\begin{methodbox}
\textbf{Photon et transition d'énergie}
\[
E=h\nu=\frac{hc}{\lambda}.
\]
Pour une transition atomique :
\[
\Delta E=E_{\text{haut}}-E_{\text{bas}}.
\]
Lors d'une émission, l'atome perd de l'énergie et émet un photon d'énergie $\Delta E$.
Lors d'une absorption, l'atome gagne de l'énergie en absorbant un photon.
\end{methodbox}

\begin{retenirbox}
\textbf{Couleurs}
\begin{itemize}
  \item Synthèse additive : addition de lumières colorées, comme sur un écran.
  \item Synthèse soustractive : absorption par des filtres ou pigments.
  \item Un objet rouge éclairé en lumière blanche apparaît rouge car il diffuse surtout le rouge.
  \item Un filtre rouge transmet surtout le rouge et absorbe les autres couleurs.
\end{itemize}
\end{retenirbox}

\begin{center}
\Large\bfseries Fin du TD
\end{center}

\end{document}